\chapter[Temperature Variations Over Isochoric Surface]{Isochoric Surfaces within the Solar Convection Zone and Temperature Variations over them}
% \section{Isochoric surfaces within the solar convection zone and temperature variations over them}
\label{ch5:appn}
%=============================================================================================================
% \lhead{\emph{Chapter 4: Appendix A}} % Write in your own chapter title to set the page header
\noindent


% \section{}
% \label{appendix}

From the basic fluid dynamical equations, we arrive at 
Equation~\ref{ch5:eq1}
involving a differentiation of $S$ with respect to $\theta$ on a surface of constant $r$.  However, the rotation of the Sun makes the Sun slightly oblate so that the solar surface is not a surface of constant $r$.  The important question is how we connect a theory involving differentiation at constant $r$ to actual observations of the solar surface.

First of all, we argue that the oblateness of the Sun is extremely small as estimated in \citet{Kuhn2012}. One can also make a rough estimate of this oblateness from theoretical considerations. If we take $\Omega$ to be constant inside the Sun for this approximate estimate, then we can introduce an effective potential
\begin{equation}
 \Phi_{\rm eff} = - \frac{G M_{\odot}}{r} - \frac{1}{2} | {\bf \Omega} \times {\bf r} |^2
 \label{ae1}
\end{equation}
inside the convection zone, which has to be constant over isobaric surfaces (see, for example, \citet{Choudhuri1998book}, Section~9.3). If $\Delta r$ is the extension of such a surface in the equatorial 
region compared to the polar region, then we can equate $\Phi_{\rm eff}$ in the polar and equatorial regions to obtain
%
\begin{equation}
 - \frac{G M_{\odot}}{r} = - \frac{G M_{\odot}}{r + \Delta r} - \frac{1}{2} \Omega^2 r^2,\nonumber
\end{equation}
from which
\begin{equation}
\Delta r = \frac{1}{2} \frac{\Omega^2 r^4}{G M_{\odot}}.
\label{ae2}
\end{equation}
Since the mass contained within the convection zone is small compared to the solar mass, we can
take $M_{\odot}$ to be the total solar mass for estimating the oblateness of isobaric surfaces
within the convection zone.  Using $\Omega/2 \pi = 420$ nHz and $r = 0.85$\Rsun, we find
\begin{equation}
\Delta r = 3.3 \; {\rm km},
\label{ae3}
\end{equation}
which is less than 0.0005\% of the solar radius. The expected oblateness of the Sun is minuscule.


We now focus our attention on isochoric surfaces over which the entropy differential would
be related with the temperature differential by Equation~\ref{ch5:eq2}.
We denote the distance measured along an isochore by $l$. It is pointed out in Appendix A of \citet{Choudhuri2021a} that we approximately have
%
\begin{equation}
\frac{dS}{dl} = \left( \frac{\pa S}{\pa \theta}\right)_r \frac{d \theta}{d l}. 
\label{ae4}
\end{equation}
Using this relation, we can put Equation~\ref{ch5:eq1} in the form
\begin{equation}
r \sin \theta \frac{\pa}{\pa z} \Omega^2 = \frac{g}{\gamma C_V} \frac{d S}{d l}, 
\label{ae5}
\end{equation}
where $d S/ dl$ is the derivative of entropy $S$ along the isochore and we have used $d l = r d \theta$
in view of the very small oblateness of the Sun.  
We can now substitute Equation~\ref{ch5:eq2} in Equation~\ref{ae5} to obtain
%
\begin{equation}
r \sin \theta \frac{\pa}{\pa z} \Omega^2 = \frac{g}{\gamma T} \frac{d}{d l} \Delta T.
\label{ae6}
\end{equation}
Now, as we move along an isochore, the change in $l$ 
is accompanied by a change in $\theta$.  We have
\begin{equation}
\frac{d}{d l} \Delta T = \frac{1}{r}
\left(\frac{d}{d \theta} \Delta T \right)_{\rm isochore}.
\label{ae7}
\end{equation}
Note that here the differenetiation is with respect to
the changing value of $\theta$ along an isochore.
On substituting Equation~\ref{ae7} into Equation~\ref{ae6}, we are readily led to Equation~\ref{ch5:eq3}. 

One fact is clear from Equation~\ref{ch5:eq3}.  If the Sun did not have a thermal wind (which would require $\Omega$
not to vary along $z$ for consistency), then the temperature would be constant over the
isochoric surface.  The isochoric surface would then be a isothermal surface and hence also an
isobaric surface.  Now,
the observed solar surface is a surface at which the optical depth becomes 1, the opacity
depending on both density and temperature. In the absence of the thermal wind, the isochoric, isothermal
and isobaric surfaces would coincide and the observed surface of the Sun would be one such surface.
%{It may be noted that in this discussion we are not considering the possibility of the very slow
%Eddington--Sweet circulation arising out of the small variation of the temperature gradient with
%latitude, since the time scale of this circulation for the Sun is larger than the age of the 
%Universe \citet{Choudhuri2021b}.
To obtain the temperature variation on the solar surface, we need to find the temperature difference
on isochoric surfaces.  For this purpose, we can use Equation~\ref{ch5:eq3}.

% \end{document}


If $\Omega (r, \theta)$ is given, then
Equation~\ref{ch5:eq3} can be integrated to obtain
$\Delta T$ along an isochore.  Because of the miniscule oblateness of the isochore, we can regard the isochore
to be a spherical surface.  However, at a conceptual level, we shall keep in mind that we would be considering
$\Delta T$ on isochoirc surfaces, since this will eventually enable us to compare with observational data
on the solar surface.

We point out one other thing, which can sometimes be
a source of confusion. By the chain rule of 
differentiation, we have
\begin{equation}
\frac{d}{dl}\Delta T = \left( \frac{\partial T}{\partial r} \right)_{\theta} \frac{d r}{dl}
+ \left( \frac{\partial T}{\partial \theta} \right)_r \frac{d \theta}{dl}. 
\label{ae8}
\end{equation}
Even though $d r/ dl$ is very small, the first term on
the right hand side of Equation~\ref{ae8} cannot be
neglected with respect to the second term because $(\partial T/ \partial r)_{\theta}$ is much larger
than $(1/r)(\partial T / \partial \theta)_r$.  As a result, we cannot write down an equation for temperature
similar to Equation~\ref{ae4} for entropy. However, when
we connect the theory with observations, the observed
temperature variations of the solar surface have to be
compared with $((d/d \theta) \Delta T)_{\rm isochore}$
appearing in Equation~\ref{ch5:eq3} and not with $(\partial T / \partial \theta)_r$ which is not of much interest to us.
